% Source: https://artofproblemsolving.com/community/c3h1277136_logic_help_math_are_champions
% Stable identity: aops-topic-1277136-q01
No. Her dad's statement has the form ``if a person is a champion, then the person is good at math.'' Emily uses the \emph{converse}, ``if a person is good at math, then the person is a champion,'' which does not follow.

In terms of sets: the statement says that the set of champions is contained in the set of people who are good at math. Emily knows only that she is in the larger set, and the larger set may contain people who are not in the smaller one. For a counterexample consistent with the dad's statement, imagine a class where the only champion is Sam, and both Sam and Emily are good at math. Then ``all champions are good at math'' is true, Emily is good at math, but Emily is not a champion. So her conclusion is not justified. (She could validly conclude ``I am not a champion'' only from ``I am not good at math''; that is the contrapositive.)
